\scanpage{193}
invariantes par $G$. On y trouve les éléments
\[
\left.\begin{aligned}
b_i&=\operatorname{Tr}x_i\\
c_i&=\det x_i
\end{aligned}\right\}\quad 1\leq i\leq4
\]
\[
t_{ij}=\deleted{\unsure{t}}t_{\{i,j\}}=\operatorname{Tr}x_i x_j
\qquad \{i,j\}\in\unsure{\mathcal P|P}_2([1,4])
\]
(ce qui fait déjà $4+4+6=14$ fonctions sur un ``espace quotient'' de dim
$4\times4-3=13$ --- il doit y avoir \deleted{\unsure{deux}}\inserted{une}
relation entre ces quantités), et enfin les
\[
t_{ijk}=\operatorname{Tr}(x_i x_j x_k)
\qquad
\deleted{1\leq i<j<k\leq4}\quad
\inserted{i,j,k\in[1,4]\text{ tous distincts}}
\]
qui sont algébriques sur l'anneau engendré par les $b_i,c_i,t_{ij}$.
Il y a $24$ fonctions $t_{ijk}$, \deleted{mais}
(\unsure{correspondants} aux \unsure{copies de Num.} à $4$ éléments) mais
\[
t_{ijk}=t_{jki}=t_{kij}
\]
donc il n'y a que $8$ fonctions distinctes, lesquelles se groupent $2$ par $2$ en
\[
(*)\quad\left\{\begin{aligned}
t_{ijk}+t_{jik}&=\Sigma_{\{i,j,k\}}
 \in\mathbb Z[b_i,c_i,t_i,b_j,c_j,t_j,b_k,c_k,t_k]\\
t_{ijk}\cdot t_{jik}&=\Pi_{\{i,j,k\}}\in\text{id}
\end{aligned}\right.
\]
\margin{192}
